\originalchapter{时间演化与跃迁率}\label{ch:golden}
\originalsection{定态叠加的演化}
若 $H$ 不显含时间, $\psi(0)=\sum_nc_n\phi_n$ 演化为 $\psi(t)=\sum_nc_n\e^{-\ii E_nt/\hbar}\phi_n$. 单一能量本征态只获得整体相位, 概率密度不变; 多能量叠加会出现相对相位和干涉. 演化算符 $U(t)=\e^{-\ii Ht/\hbar}$ 满足 $U^\dagger U=1$, 所以总概率守恒. 放射性母核的指数概率是开放通道中有效的近似描述, 不能把完整封闭体系的单位演化直接写成不守恒的态范数.

高斯波包利用自由平面波的相位逐项演化. 若初态为零平均动量、位置宽度 $a$ 的高斯, 傅里叶系数正比 $\e^{-a^2k^2}$. 因而
\begin{equation}
 \psi(x,t)=\frac{(2\pi a^2)^{-1/4}}{\sqrt{1+\ii\hbar t/(2ma^2)}}
 \exp\left[-\frac{x^2}{4a^2(1+\ii\hbar t/(2ma^2))}\right].
\end{equation}
取模平方后位置方差变成 $a^2+\hbar^2t^2/(4m^2a^2)$. 自由动量分布不变, 但不同动量分量走向不同位置, 因而波包展宽. 这把不确定关系同实际传播联系起来.

\originalsection{Heisenberg 方程与 Ehrenfest 关系}
对任意算符 $A(t)$, 从 Schrödinger 方程对 $\langle A\rangle$ 求导, 得
\begin{equation}
 \frac{\dd}{\dd t}\langle A\rangle=\frac{\ii}\hbar\langle[H,A]\rangle+\langle\partial_tA\rangle.
\end{equation}
代 $A=x,p$ 得 $\dd\langle x\rangle/\dd t=\langle p\rangle/m$ 和 $\dd\langle p\rangle/\dd t=-\langle V'(x)\rangle$. 右端一般不等于 $-V'(\langle x\rangle)$; 仅当势为至多二次式, 或波包足够窄而高阶修正小, 才与经典点粒子方程闭合. 这也是核子的平均轨迹不能替代整个量子态的原因.

\originalsection{黄金规则的时间极限}
令 $H=H_0+V$, 弱微扰从 $|i\rangle$ 激发到 $|f\rangle$. 一阶振幅为
\begin{equation}
 c_f(t)=-\frac\ii\hbar V_{fi}\int_0^t\e^{\ii(E_f-E_i)t'/\hbar}\dd t',
 \quad |c_f(t)|^2=\frac{4|V_{fi}|^2\sin^2[(E_f-E_i)t/(2\hbar)]}{(E_f-E_i)^2}.
\end{equation}
有限时间峰宽约 $\hbar/t$. 对连续而平滑的末态密度 $\rho(E)$, 积分时峰越来越窄, 面积正比 $t$. 使用
\begin{equation}\int_{-\infty}^{\infty}\frac{\sin^2(ax)}{x^2}\dd x=\pi|a|.\end{equation}, 得每单位时间的概率
\begin{equation}\label{eq:golden}
 \Gamma_{i\to f}=\frac{2\pi}{\hbar}|V_{fi}|^2\rho(E_i).
\end{equation}
条件是微扰弱、观察时间长于末态相关时间但总耗尽仍可分阶段处理、末态能谱在峰宽内近似连续. 它不能直接用于两个孤立离散能级而无限增长概率. 多个独立衰变通道的率相加, 分支比为 $\Gamma_a/\sum_b\Gamma_b$.

\originalsection{指数衰变与级联}
若每个未衰变核在短时间 $\dd t$ 的衰变概率为 $\lambda\dd t$, 宏观数目满足 $\dot N=-\lambda N$, $N=N_0\e^{-\lambda t}$. 平均寿命 $\tau=1/\lambda$, 半衰期 $t_{1/2}=\tau\log2$, 活度 $\mathcal A=\lambda N$. 把半衰期误当平均寿命会导致年龄与距离都错一个 $\log2$ 因子.

母核 $1$ 衰变到子核 $2$, $\dot N_2=\lambda_1N_1-\lambda_2N_2$. 若初始 $N_2=0$, 积分因子给
\begin{equation}
 N_2(t)=\frac{\lambda_1N_1(0)}{\lambda_2-\lambda_1}
 (\e^{-\lambda_1t}-\e^{-\lambda_2t}).
\end{equation}
当 $\lambda_1=\lambda_2=\lambda$, 极限为 $N_2=\lambda tN_1(0)\e^{-\lambda t}$. 长寿命母核与短寿命子核可趋近活度平衡, 即 $\lambda_2N_2\simeq\lambda_1N_1$; 这不意味着二者核数相同.
